% ===== §13 =====
\section{The Symptom as Fidelity to the Unsymbolizable}
\label{p26:sec:inscription}

\noindent
This section opens Part Three, which follows the persistence of drive into the life of a subject and the time of a relation. Its aim is to establish that the symptom, read through the account of the limit, is not a malfunction to be cleared away but a fidelity: a way of staying bound, in the register of repetition, to a core that was never symbolized. The section states the ordinary view of the symptom, sets against it the psychoanalytic account of repetition, gives the reading the paper proposes, and draws the consequence for reading a life.

\subsection*{The ordinary view and its limit}

If drive is the phenomenology of the limit, then its effects are not confined to a single moment but persist across the life of a subject and the life of a bond, and the first place this persistence becomes legible is the symptom. A symptom, in the psychoanalytic sense, is a formation that returns, a pattern that repeats in a subject's life with a constancy the subject does not choose and often does not want. The ordinary view treats it as a fault, a piece of the psyche gone wrong, to be corrected and dissolved, on the model of a medical symptom that points to a disorder and disappears when the disorder is cured. This view has its uses, and it has a limit, which the psychoanalytic account of repetition exposes.

\subsection*{Repetition as the return of the unsymbolized}

Freud found that the symptom does not yield to the model of a fault to be corrected. In the analysis of repetition he found a compulsion that returned upon the same, a pressure to repeat that did not serve pleasure and did not dissolve when its meaning was explained, but returned as if bound to something that explanation did not reach \parencite{freud1920beyond}. Lacan gave this its structural reading: what repeats is the subject's relation to the Real, to the point that was not taken up into the signifier, and repetition is the insistence of that unsymbolized point, its return in the only register available to it once the register of speech has failed to contain it \parencite{lacan1977xi}. The symptom, on this account, is not a fault in the symbolic but the mark of what the symbolic could not gather, returning as repetition because it could not appear as statement.

\subsection*{The symptom as fidelity}

The paper proposes to read this as a fidelity. The symptom repeats because it is bound to something that was never symbolized, a core in the Real that found no place in the signifier and so returns in the only way such a core can, as a recurrence rather than a statement. Seen this way, the symptom is the subject's way of staying bound to what could not be said. It is a fidelity held in the register of repetition, to a remainder that the register of speech could not hold. This reading does not deny the suffering a symptom can carry, and it does not counsel against the work of loosening a symptom's grip. It reframes what that work meets. What repeats is not mere error but attachment, the subject's tie to a core that matters precisely because it could not be spoken, and the constancy of the symptom is the constancy of that tie.

\begin{claim}[The symptom as fidelity]
The symptom is not only a malfunction but a fidelity: a staying-bound, in the register of repetition, to a core that was never taken up into the signifier. It returns because what it is bound to could not appear as statement and so appears as recurrence. The constancy of the symptom is the constancy of the drive's circling, made legible across a life.
\end{claim}

\subsection*{Reading a life}

Read this way, the symptom is the drive's persistence made visible in a life. What cannot be said does not therefore vanish; it presses, and its pressing takes the form of a return, a circling that shows itself over time as the recurrence of a pattern, a scene replayed, a position reassumed. The constancy of the symptom is the constancy of the drive's circling, seen across the length of a life rather than in a single turn, and its refusal to be argued away is the refusal of the Real to be talked into the symbolic. One cannot dissolve by explanation a thing whose whole nature is to be what explanation left out. The bond a person forms, the position they take up within it, the scene that recurs across their relationships, all of this carries the mark of a core that was never symbolized and goes on organizing what happens around it. The next section follows this into the time of relation, where an early signifier is shown to shape bonds that come long after.
